% ML Report for FitCoachAR (CAS 741 Extra)
% Template: ACL 2023 style, provided by Dr. Charles Welch.

\documentclass[11pt]{article}

\usepackage[]{acl}
\usepackage{times}
\usepackage{latexsym}
\usepackage[T1]{fontenc}
\usepackage[utf8]{inputenc}
\usepackage{microtype}
\usepackage{inconsolata}
\usepackage{graphicx}
\usepackage{booktabs}

\title{ML Report:\\FitCoachAR --- Pose-Based Exercise Form Feedback
with Pre-Trained Keypoint Estimators}

\author{Seyedmohamad Mirhoseininejad \\
  \texttt{mirhos5@mcmaster.ca} \\
  McMaster University --- CAS 741}

\begin{document}
\maketitle

\section{Introduction}

FitCoachAR is a real-time exercise form feedback system. It takes a
webcam stream, extracts skeletal keypoints, computes joint angles, and
emits a rep count and a form-quality signal through an AR overlay.
This report covers the ML side of the system: the pre-trained
pose-estimation backends that feed the kinematic pipeline, the
signal-processing layer on top of them, and the dynamic benchmark used
to pick between them. No model was trained from scratch for this
project, so the ML contribution is more accurately described as
applied machine learning: picking a model, evaluating it, and
integrating it with the rest of the system.

\section{Related Work}

2D multi-person keypoint estimation was popularised by
OpenPose~\citep{Cao2019OpenPose}, which introduced part-affinity
fields for associating body parts with individuals. Low-latency
on-device pose estimation then spread with the MediaPipe
framework~\citep{Lugaresi2019MediaPipe} and its BlazePose
detector~\citep{Bazarevsky2020BlazePose}, a two-stage
detector/tracker design that runs on mobile CPUs. Google's
MoveNet~\citep{MoveNet2021} is a competing lightweight model trained
on a larger pose corpus, and it is used in this project as a
comparison backend. For denoising the keypoint coordinates over time
we use a classical Kalman filter~\citep{Kalman1960}. The broader area
of recognising and evaluating human motion has been surveyed
by~\citet{Chen2021HAR} for sensor-based activity recognition.
Pose-based exercise-form feedback is a subset of that area, with the
extra constraint of real-time visual output.

\section{Dataset}

FitCoachAR does not train any model, so there is no training set in
the usual sense. The dataset used here is the evaluation dataset for
the dynamic benchmark in VnV Report \S4.1. It is a collection of
short exercise clips recorded by the author under controlled
conditions:

\begin{itemize}
  \item Subject: one adult (the author), clothed, in good lighting.
  \item Exercises: squat (primary) and bicep curl.
  \item Camera pitches: eye level and $45^{\circ}$ down.
  \item Camera orientations: front, front-$45^{\circ}$, side,
    back-$45^{\circ}$, back.
  \item Format: H.264 MP4, 30~FPS, 720p.
\end{itemize}

No manual per-frame annotation was done. Ground truth for joint angles
comes from two sources. The first is a closed-form analytical dataset
of five synthetic three-point configurations with known angles
spanning $[0^\circ, 180^\circ]$; this is used in T8 of the VnV Report.
The second is the rep-counting ground truth, which is a physical
property of the video (``was a valid rep performed or not?'') and can
be labelled by watching the video. So the ML task reported here is
not supervised learning over a labelled dataset. It is model selection
and evaluation using a stability proxy. This is made explicit in the
Results section.

\section{Features}

The features the kinematic engine consumes are the 3D coordinates of
six body landmarks: hip, knee and ankle for the lower body, and
shoulder, elbow and wrist for the upper body. These come from the
pre-trained pose backend and are indexed into its landmark list using
the backend's published schema. No learned embedding is applied on
top. Keeping the feature pipeline light is a deliberate choice: it
keeps end-to-end latency inside the real-time budget and within the
NFR1 $\pm 5^\circ$ accuracy target. Temporal smoothing is done per
axis using a Kalman filter~\citep{Kalman1960}, implemented in module
M8 (\texttt{KalmanLandmarkSmoother}); this cuts down frame-to-frame
jitter before the angles are computed.

\section{Implementation}

\subsection{Pose estimation backends}

Three backend configurations were evaluated:

\begin{itemize}
  \item MediaPipe 2D, the default production backend. Uses
    BlazePose~\citep{Bazarevsky2020BlazePose} via the MediaPipe Python
    package in 2D mode, producing 33 landmarks per frame with $(x,y)$
    coordinates normalised to the image and a per-landmark visibility
    score.
  \item MediaPipe 3D, the same BlazePose model with ``world
    landmarks'' enabled. This adds a $z$ coordinate in metres
    relative to the subject's hip. $z$ comes from a learned
    regression head.
  \item MoveNet 3D~\citep{MoveNet2021}, Google's lightweight model,
    used via its TensorFlow-Hub distribution. Produces 17 keypoints.
\end{itemize}

The three backends are treated as black boxes. Module M4
(\texttt{m4\_pose\_tracking.py}) wraps them behind a common
\texttt{IPoseBackend} interface so the rest of the pipeline is
backend-agnostic.

\subsection{Downstream analytical pipeline}

Once the backend produces landmarks, the analytical layer computes
joint angles using the vector dot-product formulation
\begin{equation}
  \theta = \arccos\!\left(\frac{\overrightarrow{ba}\cdot
  \overrightarrow{bc}}{\|\overrightarrow{ba}\|\;\|\overrightarrow{bc}\|}\right)
  \cdot \frac{180}{\pi}
  \label{eq:angle}
\end{equation}
with a clamp on the cosine argument to avoid floating-point overshoot.
These angles feed the four-state exercise FSM in M6, which counts
valid reps with a minimum-duration guard $t_{\min}$. The FSM is plain
control logic and has no learned parameters. It is included here
because it consumes the ML-derived features.

\subsection{Training procedure}

None. All three pose models are used at their published checkpoints
with no fine-tuning. The evaluation dataset is one subject under
controlled conditions, which is too small and too biased to justify
fine-tuning. And the quality that matters for this application is how
well the backend generalises, not how well it fits the author's
silhouette.

\section{Results and Evaluation}

\subsection{Evaluation protocol}

There is no dense per-frame angle ground truth for the recorded
dataset, so two complementary metrics are used instead.

\begin{enumerate}
  \item Analytical MAE (geometric). For the five synthetic angles in
    test T8 (VnV Report \S6.9), the mean absolute error of the
    computed angle against the closed-form ground truth is below
    $0.01^{\circ}$ in every case. This confirms the kinematic engine
    itself is correct, so any angular error on real video has to come
    from the pose backend.
  \item Stability proxy $\overline{|\Delta\theta|}$. For the recorded
    dataset we compute
    $\overline{|\Delta\theta|} = \frac{1}{N{-}1}\sum_{i=1}^{N-1}
    |\theta_i - \theta_{i-1}|$ per backend. This is a stability
    measure, not an error measure: lower values mean a smoother
    output signal over time when the true angle is changing smoothly.
\end{enumerate}

\subsection{Offline benchmark}

Table~\ref{tab:benchmark} reports the dynamic benchmark results from
the VnV Report (\S4.1), carried over here so the ML Report is
self-contained.

\begin{table}[h]
  \centering
  \begin{tabular}{lccc}
    \toprule
    \textbf{Backend} & \textbf{Latency} & \textbf{Knee $\overline{|\Delta\theta|}$} &
    \textbf{Elbow $\overline{|\Delta\theta|}$} \\
    & (ms) & ($^\circ$) & ($^\circ$) \\
    \midrule
    MediaPipe 2D & ${\sim}39$ & ${\sim}0.7$ & ${\sim}3.8$ \\
    MediaPipe 3D & ${\sim}41$ & ${\sim}1.2$ & ${\sim}2.3$ \\
    MoveNet 3D   & ${\sim}46$ & ${\sim}1.4$ & ${\sim}5.4$ \\
    \bottomrule
  \end{tabular}
  \caption{Offline benchmark across the three pose backends.
  $\overline{|\Delta\theta|}$ is mean absolute frame-to-frame angle
  delta (lower~=~more stable); it is \emph{not} error vs.\ ground
  truth.}
  \label{tab:benchmark}
\end{table}

\subsection{Analysis}

MediaPipe 2D wins on knee-angle stability and on latency, which is
the reason it is the production backend. Two patterns in the table
are worth calling out:

\begin{itemize}
  \item MediaPipe 3D is better than 2D on the elbow but worse on the
    knee. The $z$-regression head seems to help on distal joints like
    the elbow, where in-plane projection ambiguity is largest, but it
    adds noise on large-range-of-motion joints like the knee. This
    matches the known observation that BlazePose's $z$-head is only
    approximate and is mostly informative near frontal poses.
  \item MoveNet is the worst on both joints. It is a 17-keypoint
    model, which is coarser than BlazePose's 33-landmark schema. The
    noisier output, combined with the same Kalman smoother, gives a
    worse end-to-end stability result.
\end{itemize}

\subsection{Limitations and failure cases}

The benchmark has some real limitations:

\begin{itemize}
  \item Single-subject evaluation. All clips are one adult body type
    in controlled lighting. How well the backends generalise to other
    body types, clothing and lighting is inherited from the pose
    backend and not measured directly here.
  \item Stability is not correctness. A backend that is consistently
    wrong by $10^{\circ}$ would still score well on
    $\overline{|\Delta\theta|}$. The T8 analytical test partly
    compensates by verifying the kinematic engine, but it does not
    verify the pose backend itself against a physical goniometer. An
    end-to-end goniometer validation was planned but not done. It is
    still open.
  \item No occlusion sweep. T9 in the VnV Plan was designed but
    deferred, so behaviour under severe limb occlusion is not
    quantified.
\end{itemize}

\section{Feedback and Plans}

Two peer-review comments during the term touched the ML side. The
first was about the term ``MAE''. The earlier version of this report
and of the VnV Report used MAE for the stability proxy, which
conflated it with the actual Mean Absolute Error reported in T8. The
reviewer flagged this, and both documents now use
$\overline{|\Delta\theta|}$ for the stability metric (see VnV Report
\S4.1). The second comment questioned the decision to treat the pose
backend as a black box and to put effort into the downstream
analytical layer. The benchmark numbers above are the answer: most of
the end-to-end noise comes from the backend itself, and M8 smoothing
cannot recover information the backend has already lost. Future work
would run a goniometer validation on a sample of frames, which would
turn $\overline{|\Delta\theta|}$ into a real MAE; quantify occlusion
behaviour via T9; and broaden the evaluation set to more subjects and
lighting conditions.

\section*{Ethical and Data Considerations}

The evaluation dataset was produced by and of the author, so no
third-party consent was required. For a real deployment, two ethical
risks are carried over from the pre-trained backends and should be
stated explicitly:

\begin{itemize}
  \item Demographic bias. Pose-estimation accuracy is known to vary
    with skin tone, body type, clothing and lighting. A form-feedback
    system that is less accurate on some users than others will give
    worse coaching to those users. This is a fairness concern.
  \item Privacy. The system runs entirely on the user's browser and a
    user-controlled backend. No video leaves the client unless the
    user opts in to server-side processing. The Django WebSocket
    endpoint does not store frames. This is an architectural
    decision, but it is worth mentioning here because pose data is
    biometric.
\end{itemize}

\bibliography{custom}

\end{document}
